\section{Limitations}\label{sec:limitations}

\paragraph{Simulation.} The simulation study uses subnetworks of 400 and 800 genes, a single network score threshold and a simple patient model with a shared component and independent errors. The shared component cancels from the plain change, and the difference of two independent gamma errors is symmetric, so the model does not test correlated changes along pathways, heterogeneous changes between patients or tissue composition. The Level~2 test was not run under a complete null with the final reference. The matched experiment plants one to three centres of a single kind and two effect sizes; it exchanges a centre with a neighbour, which changes more than the presence of a hole; and it does not record whether the planted class survives in the full filtration. Whether $\Hone$ adds information to the baselines jointly was not tested. The network-adjusted score used $R=19$ relabellings per patient in the simulations, fewer than the 199 used for the real data, and its simulated reference was stratified by degree and clustering only. The null calibration with normal and outlier noise and the matched experiment were repeated on reserved verification seeds with the final code; the heavy-tailed and skewed null runs and the hub-change run were made on development seeds only.

\paragraph{Inference.} Validity of the tests is validity under $H_{\mathrm{sym}}$, the exchangeability of each patient's two specimens. In real data this can fail for technical reasons, such as tumour and normal specimens processed on different plates, and adjacent normal tissue is not healthy tissue: it can carry field effects, and both specimens are mixtures of cell types. A rejection therefore shows a paired difference in directional topology, not its cause. For the gene family, family-wise error control is established under the complete null; under a partial null, the validity of each unaffected gene's $p$-value conditionally on a family selected from all genes has not been shown (Section~\ref{sec:inference}). Cycle representatives are deterministic but not canonical; gene scores depend on them, and we checked only the sign of each score under one alternative tie rule.

\paragraph{Network and data.} STRING edges include indirect associations such as shared complex membership, and filling 3-cliques is a modelling choice; cycles describe arrangements of assay values on a reference graph. The methylation value of a gene is a median of promoter-probe $M$-values, and statements about hypermethylation concern the analysed promoter-associated genes, not genome-wide methylation. The composition adjustment uses RNA-derived stromal and immune proxies only, applies them to the methylation layer as well, and is exploratory at Level~2. The alternative network thresholds (STRING scores 400 and 900) were applied to the genes of the primary network rather than to each network's own gene set, because the preprocessed data hold only those genes.

\paragraph{Cohort.} The TCGA-BRCA cohort has \cohortN{} eligible pairs, a single source and a cross-sectional design. Independent confirmation is planned in the Hong Kong breast cancer cohort (dbGaP phs001870, controlled access; paired tumour and adjacent normal tissue with RNA sequencing and EPIC methylation profiling), with network, gene universe, preprocessing and code frozen before those data are opened; the number of patients with both layers in both tissues will be reported once access is granted. A data access request has been drafted.
